\section{Research Approach}

The current research approach begins from a practical observation. Existing editing methods do not fail for the same reason. Some fail because they are not stable enough under long sequential use, while others remain strong only by paying a high computational or operational cost. The thesis therefore does not assume that lifelong editing should be solved by searching for a single universally dominant editor. Instead, it studies the conditions under which different editing mechanisms are suitable and uses those observations to motivate a hybrid design.

At the current stage, the work is organized in two layers. The first layer is an empirical layer, in which representative baselines are implemented and evaluated under the same sequential protocol. The purpose of this layer is to expose the actual trade-off between stability, locality, portability, and editing cost on a common model and benchmark setting. The second layer is a design layer, in which the empirical findings are used to motivate a complementary lifelong editing architecture whose fast and slow behaviors are not forced into one mechanism prematurely.
